\section{Introduction}\label{sec:intro}
Internet of Things (IoT) systems connect physical devices that carry sensors and network interfaces and that exchange data with remote servers~\cite{atzori2010,gubbi2013}. These systems now support agriculture, health monitoring, smart homes, and manufacturing. As deployments grow, operators must detect anomalies that indicate physical faults or cyberattacks, because undetected anomalies cause unplanned downtime, quality loss, and security incidents.

Anomalies in IoT systems originate from physical faults and from cyberattacks. A worn bearing or an unbalanced motor changes the vibration that attached sensors record, and an unnoticed fault can halt a machine without warning~\cite{wang2018,gummadi2024xaiiot}. In networks, consumer devices such as doorbells, thermostats, baby monitors, and security cameras can be infected by botnets such as Mirai and Gafgyt and then used to attack other hosts~\cite{meidan2018nbaiot,antonakakis2017}. In both cases, the observed readings or traffic deviate from normal behavior~\cite{chandola2009}.

Early detection of such anomalies supports predictive maintenance in manufacturing, since repairs can be planned before a failure occurs~\cite{wang2018}. Low-cost MEMS accelerometers cost a fraction of the price of piezoelectric sensors but sample at low rates, which yields sparse data~\cite{gummadi2024xaiiot}. In networks, signature-based detection identifies known attacks but struggles with new threats and changing attack patterns~\cite{asad2022}. Learned anomaly detectors can close part of this gap, provided that their models respect the memory and latency limits of IoT gateways.

Prior work has applied machine learning to anomaly detection in IoT sensor data and network traffic~\cite{vinayakumar2019,tama2019,bibers2024,gummadi2024xaiiot}. However, most of the existing work selects and tunes models manually, and the conclusions depend on the dataset and on the evaluation protocol. For example, a neural network performed best on sensor data in~\cite{gummadi2024xaiiot}, whereas ensemble methods performed best on botnet traffic in the same study.

Automated machine learning (AutoML) frameworks aim to remove this manual effort by searching jointly over learning algorithms and their hyperparameters~\cite{thornton2013,feurer2015,ledell2020}. Automated pipelines have also been proposed for intrusion detection~\cite{yang2024automl,injadat2021,guntur2025autoai} and for industrial IoT traffic~\cite{sezgin2023}. However, these studies address network traffic, and independent comparisons between frameworks remain rare.

On the contrary, evaluating AutoML on IoT data requires particular care. Public security datasets are known to contain labeling errors, duplicated records, and leakage between training and test partitions~\cite{engelen2021,lanvin2023,kapoor2023}. IoT datasets may share these problems, and then the reported advantage of a framework may reflect the evaluation protocol instead of the framework itself.

There are articles discussing evaluation pitfalls in security and machine learning~\cite{arp2022,kapoor2023}, none of which however examines how these pitfalls change conclusions about AutoML frameworks on IoT data. Our work bridges this gap by auditing public IoT datasets, constructing leakage-aware protocols, and comparing AutoML frameworks under identical conditions.

There are recent works~\cite{guntur2025autoai,velnambi2026} that have compared Auto-sklearn and H2O AutoML on intrusion detection data. However, these works have the following limitations that our study addresses. They address network traffic and do not consider IoT sensor data. In our earlier study~\cite{velnambi2026}, the frameworks received different training data and search budgets, so the observed differences could not be attributed to the frameworks. Moreover, neither work audits the benchmark data before it ranks the methods. The difference between AutoML-IoT and previous related work is shown in Table~\ref{tab:related}.
\begin{table*}[t]
\centering
\caption{Comparison between prior work and AutoML-IoT in terms of the available features (to our reading of the cited works).}
\label{tab:related}
\scriptsize
\setlength{\tabcolsep}{3pt}
\begin{tabular}{@{}p{1.95in}ccccccc@{}}
\toprule
System & IoT data & Automated selection & Independent comparison & Data audit & Leakage-aware protocol & Cost measured & Explanation \\
\midrule
Auto-WEKA, TPOT, Auto-sklearn, H2O AutoML, AutoGluon-Tabular \cite{thornton2013,olson2016,feurer2015,ledell2020,erickson2020} & \xmark & \cmark & \xmark & \xmark & \xmark & \xmark & \xmark \\
AMLB benchmark \cite{gijsbers2022} & \xmark & \cmark & \cmark & \xmark & \xmark & \cmark & \xmark \\
XAI-IoT \cite{gummadi2024xaiiot} & \cmark & \xmark & \xmark & \xmark & \xmark & \xmark & \cmark \\
AutoAI-IDS \cite{guntur2025autoai} & \xmark & \cmark & \cmark & \xmark & \xmark & \cmark & \xmark \\
Earlier AutoML study \cite{velnambi2026} & \xmark & \cmark & \cmark & \xmark & \xmark & \cmark & \xmark \\
CICIDS2017 audits \cite{engelen2021,lanvin2023} & \xmark & \xmark & \xmark & \cmark & \xmark & \xmark & \xmark \\
Pitfall studies \cite{arp2022,kapoor2023} & \xmark & \xmark & \xmark & \cmark & \cmark & \xmark & \xmark \\
\midrule
\textbf{AutoML-IoT (this work)} & \cmark & \cmark & \cmark & \cmark & \cmark & \cmark & \cmark \\
\bottomrule
\end{tabular}
\end{table*}


\textbf{Our contributions:} In this paper, we audit two public IoT datasets, a vibration dataset collected on a motor testbed and a traffic dataset collected from commercial devices. We then design an evaluation pipeline that we call AutoML-IoT. The pipeline consists of consecutive steps for data audit, protocol construction, controlled search, measurement, explanation, and statistical comparison, as shown in Figure~\ref{fig:pipeline}. In every step, both frameworks receive identical inputs, and default baseline models serve as a reference.
\begin{figure*}[t]
\centering
\begin{tikzpicture}[
  node distance=0.5cm and 0.45cm,
  box/.style={draw=black!65, rounded corners=3pt, fill=black!5, text width=2.12in, align=left, inner sep=4pt, font=\scriptsize},
  arr/.style={-{Stealth[length=5pt]}, thick, draw=black!70}
]
\node[box] (s1) {\textbf{Step 1: Data audit.} Duplicates, label-ambiguous vectors, accuracy ceiling, and label blocks as inferred sessions. Output: raw and cleaned data.};
\node[box, right=of s1] (s2) {\textbf{Step 2: Protocol construction.} Random and block splits, raw and cleaned data, training distribution, and fixed seeds. Output: identical training and test arrays.};
\node[box, right=of s2] (s3) {\textbf{Step 3: Controlled search.} Auto-sklearn, H2O AutoML, and default baselines under identical data, budget, and threads. Output: selected models.};
\node[box, below=0.75cm of s3] (s4) {\textbf{Step 4: Measurement.} Macro-averaged quality metrics, per-class recall, memory, model size, and latency. Output: per-run records.};
\node[box, left=of s4] (s5) {\textbf{Step 5: Explanation.} Permutation importance of the selected models and agreement between frameworks. Output: feature rankings.};
\node[box, left=of s5] (s6) {\textbf{Step 6: Paired comparison.} Differences by seed, summary tables, and recommendations. Output: guidance for model selection.};
\draw[arr] (s1) -- (s2);
\draw[arr] (s2) -- (s3);
\draw[arr] (s3) -- (s4);
\draw[arr] (s4) -- (s5);
\draw[arr] (s5) -- (s6);
\end{tikzpicture}
\caption{A high-level overview of AutoML-IoT. A labeled IoT dataset passes through six consecutive steps that produce comparable records for every method.}
\label{fig:pipeline}
\end{figure*}


We show that the evaluation protocol changes the reported performance more than the choice of framework. In particular, holding out entire recording blocks lowers the macro-averaged F1 score of all methods, whereas the differences between the frameworks stay within the variation across random seeds. We also show that label-ambiguous records bound the attainable accuracy on the traffic data and that both frameworks perform close to a default random forest. The frameworks differ mainly in the size and latency of the selected models. In obtaining these results, we address challenges that are specific to IoT data. These include the absence of session identifiers in the sensor data, repeated and conflicting records in the traffic data, and different validation strategies in the frameworks.

In summary, this paper makes the following contributions:
\begin{enumerate}
\item[(1)] We propose \emph{AutoML-IoT}, a controlled evaluation pipeline for AutoML frameworks on IoT anomaly detection. The pipeline covers data audit, leakage-aware protocol construction, controlled search, measurement of quality and cost, explanation, and paired statistical comparison.
\item[(2)] We audit a sensor dataset and a traffic dataset and quantify duplication, label ambiguity, and temporal block structure. For the traffic data, we derive an accuracy ceiling that follows from identical inputs with different labels.
\item[(3)] We compare Auto-sklearn and H2O AutoML under identical data, budgets, thread counts, and seeds, and we add default baseline models. We analyze split strategy, data cleaning, and class balancing as separate factors.
\item[(4)] We report macro-averaged quality metrics, per-class recall, selected models, memory, model size, and prediction latency for both frameworks in complete tables.
\item[(5)] We explain the selected models with permutation importance and relate the result to hand-built models from earlier work.
\item[(6)] We derive recommendations for evaluating AutoML on IoT data, and we release code and per-run results.
\end{enumerate}
